\documentclass[11pt,a4paper]{article}
\usepackage[margin=25mm]{geometry}
\usepackage{amsmath,amssymb,amsthm,hyperref}
\setlength{\emergencystretch}{2em}
\newtheorem{lemma}{Lemma}
\newtheorem{theorem}[lemma]{Theorem}
\newcommand{\gen}[1]{\langle#1\rangle}
\title{Support bounds without a parity hypothesis}
\author{Human Compression project: a component of the L019 recovery}
\date{}
\begin{document}\maketitle
Let $G=\gen{\rho_0,\ldots,\rho_{r-1}}$ be a transitive string
action by involutions, with connected noncommutation diagram and
all maximal omissions intransitive. Let a unique eligible
$i$-edge join $a\in A$ to $b\in B$, the two omission orbits.
Assume $|A|,|B|\ge2$ and $L=G_{<i}$ fixes $B$ pointwise.
Write $R=G_{>i}$. These two groups commute, their product is
transitive on $A$, and $R$ is transitive on $B$.

We acquire the exact unique-crossing and fixed-prefix arguments
from the earlier work. In particular, if $l>i$ acts on either
side, there is a simple path from that side's crossing endpoint
to an $l$-edge, using only earlier indices different from $i$.
The path with its last edge has at least $l-i$ edges; its prefix,
excluding that last edge, is fixed pointwise by $G_{>l}$.
Every intermediate index $i+1,\ldots,l-1$ occurs. These assertions
are generic and use no evenness or intersection property.

\begin{lemma}[The commuting-orbit counts]\label{orbits}
Suppose $i>0$ and $L$ is intransitive on $A$. Its orbits in
$A$ have a common size $k\ge i+1$. Put $\beta=a^L$ and
$\tau=\rho_{i+1}$. Then $\tau$ exchanges $\beta$ with a
distinct orbit $\beta'$. All marks of index at least $i+2$
fix $\beta$ pointwise, and those of index at least $i+3$
fix $\beta'$ pointwise.
If $\rho_h$ acts nontrivially on $A$, where $d=h-i\ge2$,
there are at least $d$ such orbits. If the first active orbit
reached by the earlier marks is moved to another orbit by
$\rho_h$, there are at least $d+1$.
\end{lemma}
\begin{proof}
The group $R$ transports the $L$-orbits equivariantly and
transitively, so their sizes agree. A subfamily of the $L$
marks transitive on one orbit is transitive on all of them.
Together with $R$ it is transitive on both $A,B$. The acquired
minimal-transitivity lemma therefore requires every $L$ mark
in this subfamily. Its orbit forest gives $i\le k-1$.

All indices at least $i+2$ fix $a$ by the unique crossing
and commute with $L$, hence fix $\beta$ pointwise. Therefore
$\tau$ must send $\beta$ elsewhere, or it would be invariant
under $G_i$. This gives $\beta'$. Commutation with $\tau$
then makes indices at least $i+3$ fix $\beta'$ pointwise.

Consider the $L$-orbit graph using indices strictly between
$i$ and $h$, starting at $\beta$. Some reachable orbit has
nontrivial $\rho_h$-action. Otherwise $\rho_h$ would fix
their union pointwise; every higher mark does so too, since
it fixes $\beta$ and commutes with all the path marks. Their
union would be $G_i$-invariant, hence all of $A$, contradicting
the assumed action of $\rho_h$.
Choose a shortest orbit path to such an orbit. Every intermediate
index occurs in its word: if one were missing, commute the
higher-index letters to the beginning. They fix $a$, and
the remaining lower-index word commutes with $\rho_h$ and
starts at a point fixed by it. This would make its endpoint
fixed by $\rho_h$. At an active invariant orbit this is
impossible, since a nonidentity centralizer of a transitive
$L$-action has no fixed point; if that orbit is moved, it
is impossible immediately. The path has at least $d-1$
edges and $d$ distinct orbits. Every earlier orbit is fixed
pointwise by $\rho_h$, by minimality. If the last one is
moved, its image is a further orbit, proving the last count.
\end{proof}

\begin{lemma}[An internal action at distance two]\label{internal}
Under Lemma~\ref{orbits}, suppose $\sigma=\rho_{i+2}$
preserves $\beta'$ and is nontrivial there. Then
\[
              A=\beta\cup\beta',\qquad k\ge4,
              \qquad i\le k/2-1.
\]
\end{lemma}
\begin{proof}
Every later generator fixes $\beta\cup\beta'$ pointwise;
$L,\tau,\sigma$ preserve its union. Transitivity on $A$
therefore gives the asserted union. On $\beta'$, $\sigma$
centralizes the transitive $L$-action, so is fixed-point-free.
Its $k/2$ pairs are $L$-invariant, and $\sigma^\tau$ supplies
the corresponding pairs in $\beta$.
Choose a subfamily $M$ of the $L$ marks transitive on these
pairs, with $|M|\le k/2-1$. The group $\gen{M,\sigma^\tau}$
is transitive on $\beta$, and $\gen{M,\sigma}$ on $\beta'$.
Together with $\tau$ they are transitive on $A$. The whole
$R$ remains transitive on $B$. Minimal transitivity at the
crossing forces $M$ to contain all $L$ marks. This proves
$i\le k/2-1$; since $i>0$, $k\ge4$.
\end{proof}

\begin{lemma}[A transitive first side]\label{transitive}
If $i>0$, $L$ is transitive on $A$, and a mark after $i$
acts nontrivially on $A$, that mark has index $i+1$.
Moreover $|A|=2m$ with $i\le m-1$.
\end{lemma}
\begin{proof}
Every mark of index at least $i+2$ fixes $a$ and centralizes
the transitive $L$-action, so fixes $A$ pointwise. Thus the
nontrivial action belongs to $\tau=\rho_{i+1}$, whose
restriction is fixed-point-free. Its $m$ pairs are $L$-invariant.
Choose at most $m-1$ of the $L$ marks transitive on the pairs.
Together with $\tau$ these are transitive on $A$, and the
whole $R$ is transitive on $B$. Minimal transitivity forces
all $i$ marks to be present, as claimed.
\end{proof}

\begin{lemma}[The outside side alone]\label{outside}
If $h>i$ and $\rho_h$ acts nontrivially on $A$, then
\[
                         |A|\ge i+h+1.
\]
For $i=0$ no hypothesis that $G_{<i}$ be nontrivial is needed.
\end{lemma}
\begin{proof}
For $i=0$ the acquired path has at least $h+1$ vertices in
$A$. For $i>0$ and transitive $L$, Lemma~\ref{transitive}
gives $h=i+1$ and $|A|\ge2i+2$. Otherwise use
Lemma~\ref{orbits} and set $d=h-i$.
For $d=1$, $|A|\ge2k\ge2i+2$. For $d=2$, the first
active orbit is $\beta'$. If it is moved, $|A|\ge3k\ge3i+3$
is sufficient. If it is preserved nontrivially, Lemma~\ref{internal}
gives $|A|=2k\ge4i+4$. It cannot be fixed pointwise, since
then $\beta\cup\beta'$ would be invariant under every mark
of $G_i$, and $\rho_h$ would be trivial on $A$.
For $d\ge3$,
\[
 |A|\ge d(i+1)\ge 2i+d+1=i+h+1,
\]
where $i(d-2)\ge1$. This includes the possible equality
$i=1,d=3$; no parity argument is used to discard it.
\end{proof}

\begin{lemma}[A last mark acting on the first side]\label{last}
If $i>0$ and $\rho_{r-1}$ acts nontrivially on $A$, then
$n\ge2r$. If $i=0$ and the last mark acts on both sides,
the same conclusion holds.
\end{lemma}
\begin{proof}
For $i=0$, take the two disjoint acquired paths, each with
$r$ vertices. Now suppose $i>0$. If $L$ is transitive on
$A$, Lemma~\ref{transitive} gives $r-1=i+1$. On $B$, $R$
is the single involution $\tau$, so $|B|=2$. Thus
$n=2m+2\ge2i+4=2r$.

Suppose $L$ is intransitive. First let the last mark fix
$B$ pointwise. Its image on the $L$-orbits in $A$ cannot
belong to the group generated by the other $R$ images:
otherwise those images would remain transitive on the orbit
set, and $L$ with the remaining $R$ marks would be transitive
on both $A,B$. This contradicts minimal transitivity.
In particular its block image is nontrivial. Put
$d=r-1-i\ge2$, since $\tau$ moves $b$. A shortest block
path from $\beta$ to a block moved by the last mark uses
each intermediate index, by the same no-gap argument as
Lemma~\ref{orbits}. It visits at least $d$ distinct blocks;
the last mark sends its final block to one more. Hence
\[
 n\ge(i+1)(d+1)+2=2r+(i-1)(d-1)\ge2r.
\]

Finally suppose the last mark acts on both sides. If $d=1$,
$A=\beta\cup\beta'$ and $|B|=2$, giving
$n=2k+2\ge2i+4=2r$. If $d=2$, the active orbit $\beta'$
is either moved, giving $|A|\ge3k$, or preserved nontrivially,
giving Lemma~\ref{internal}. In the first case the path in
$B$ gives $|B|\ge3$, so $n\ge3i+6\ge2r$. In the second,
$n\ge2k+3$ and $r=i+3\le k/2+2$, which again imply
$n\ge2r$ because $k\ge4$. For $d\ge3$, the orbit count
and the path in $B$ give
\[
 n\ge d(k+1)+1\ge d(i+2)+1
                =2r+i(d-2)-1\ge2r.
\]
These cases exhaust the last mark's possible action.
\end{proof}

\begin{theorem}[The head support bound]\label{head}
Assume $2r>n+2$ and at least one mark after $i$ is shared
by both sides. Let $h>i$ be the largest shared index.
Then $h<r-1$. If $i=0$, interchange $A,B$ when necessary
so that the final tail acts on $B$. For $i>0$ this is
already its side. With $X=\operatorname{supp}(G_{>h})$,
\[
             X\subseteq B\setminus\{b\},\qquad
             |\operatorname{Fix}(G_{>h})|\ge2h+1.
\]
\end{theorem}
\begin{proof}
Lemma~\ref{last} excludes $h=r-1$. Every later mark is
nonidentity and acts on exactly one side. Connectedness
of the noncommutation diagram forces the same side for
all of them: a change of sides would make adjacent marks
commute. For $i>0$, Lemma~\ref{last} excludes $A$ as this
side. The tail therefore fixes all of $A$ and also fixes
$b$ by the unique crossing.

For $i=0$, the path in $A$ including its last $h$-edge has
at least $h+1$ vertices, all fixed by the tail. The path
prefix in $B$ has at least $h$ vertices, also fixed. These
disjoint sets give $2h+1$ fixed points.

Let $i>0$. If $L$ is transitive, then $h=i+1$ and
$|A|+1\ge2i+3=2h+1$ by Lemma~\ref{transitive}.
Otherwise put $d=h-i$. The fixed prefix in $B$ has at
least $d$ vertices, so the fixed set has size at least
$|A|+d$. If $d=1$, this is at least $2k+1\ge2h+1$.
If $d=2$ and the active orbit is moved, it is at least
$3k+2\ge3i+5\ge2h+1$. In the internal case it is at
least $2k+2$, whereas $i\le k/2-1$ and $k\ge4$ imply
$2h+1=2i+5\le k+3\le2k+2$.
For $d\ge3$, it is at least
\[
 d(k+1)\ge d(i+2)=2h+1+i(d-2)-1\ge2h+1.
\]
The support inclusion follows from the already proved
fixation of $A$ and $b$.
\end{proof}

\paragraph{Attribution and boundary.}
The source is Cameron--Fernandes--Leemans,
\emph{The number of string C-groups of high rank},
arXiv:2212.12723v2, Propositions~5.8--5.12. Its path and
commuting-orbit strategy, and the complete generic path proofs
already acquired for L044, remain attributed. The preceding
L044 orbit-distance, transitive-side and internal-block arguments
are reworked here with minimal transitivity replacing faithful
quotient rank bounds. Every change in hypothesis is proved:
evenness and the alternating small-degree assumptions are not
used, and the old parity exclusions of equality are not claimed.
The support $X$ is on the side \emph{moved} by the final tail;
the source's contrary orbit wording is not retained.

This manuscript establishes the exact weak inequalities needed
in the strict L019 range. It does not claim the stronger equality
classifications of the source or the alternating manuscripts,
or yet establish existence of a perfect split. No new census,
whole-original, conjecture, Lean or measurement credit is asserted.
\end{document}
